\documentclass[10pt,a4paper]{amsart}
\usepackage{amsmath,amssymb,amsthm,graphicx}
% For commutative diagrams we use Xy-pic
\usepackage[all]{xy}

% We use 2cell for 2-commutative diagrams.
\xyoption{2cell}
% Theorem environments.
%
\theoremstyle{plain}
\newtheorem{theorem}[subsection]{Theorem}
\newtheorem{proposition}[subsection]{Proposition}
\newtheorem{lemma}[subsection]{Lemma}

\theoremstyle{definition}
\newtheorem{definition}[subsection]{Definition}
\newtheorem{example}[subsection]{Example}
\newtheorem{exercise}[subsection]{Exercise}
\newtheorem{situation}[subsection]{Situation}

\theoremstyle{remark}
\newtheorem{remark}[subsection]{Remark}
\newtheorem{remarks}[subsection]{Remarks}

\numberwithin{equation}{subsection}

% Macros
%
\def\lim{\mathop{\mathrm{lim}}\nolimits}
\def\colim{\mathop{\mathrm{colim}}\nolimits}
\def\Spec{\mathop{\mathrm{Spec}}}
\def\Hom{\mathop{\mathrm{Hom}}\nolimits}
\def\Ext{\mathop{\mathrm{Ext}}\nolimits}
\def\SheafHom{\mathop{\mathcal{H}\!\mathit{om}}\nolimits}
\def\SheafExt{\mathop{\mathcal{E}\!\mathit{xt}}\nolimits}
\def\Sch{\mathit{Sch}}
\def\Mor{\mathop{\mathrm{Mor}}\nolimits}
\def\Ob{\mathop{\mathrm{Ob}}\nolimits}
\def\Sh{\mathop{\mathit{Sh}}\nolimits}
\def\NL{\mathop{N\!L}\nolimits}
\def\CH{\mathop{\mathrm{CH}}\nolimits}
\def\proetale{{pro\text{-}\acute{e}tale}}
\def\etale{{\acute{e}tale}}
\def\QCoh{\mathit{QCoh}}
\def\Ker{\mathop{\mathrm{Ker}}}
\def\Im{\mathop{\mathrm{Im}}}
\def\Coker{\mathop{\mathrm{Coker}}}
\def\Coim{\mathop{\mathrm{Coim}}}

% Boxtimes
%
\DeclareMathSymbol{\boxtimes}{\mathbin}{AMSa}{"02}

%
% Macros for moduli stacks/spaces
%
\def\QCohstack{\mathcal{QC}\!\mathit{oh}}
\def\Cohstack{\mathcal{C}\!\mathit{oh}}
\def\Spacesstack{\mathcal{S}\!\mathit{paces}}
\def\Quotfunctor{\mathrm{Quot}}
\def\Hilbfunctor{\mathrm{Hilb}}
\def\Curvesstack{\mathcal{C}\!\mathit{urves}}
\def\Polarizedstack{\mathcal{P}\!\mathit{olarized}}
\def\Complexesstack{\mathcal{C}\!\mathit{omplexes}}
% \Pic is the operator that assigns to X its picard group, usage \Pic(X)
% \Picardstack_{X/B} denotes the Picard stack of X over B
% \Picardfunctor_{X/B} denotes the Picard functor of X over B
\def\Pic{\mathop{\mathrm{Pic}}\nolimits}
\def\Picardstack{\mathcal{P}\!\mathit{ic}}
\def\Picardfunctor{\mathrm{Pic}}
\def\Deformationcategory{\mathcal{D}\!\mathit{ef}}
\usepackage{cleveref}
\crefname{originalreference}{original reference}{original references}
\Crefname{originalreference}{Original reference}{Original references}
\makeatletter
\newlabel{definition-locally-acyclic}{{0GJP}{source}{}{sourcepointer-1}{}}
\newlabel{definition-locally-acyclic@cref}{{[originalreference][0][]0GJP}{[1][0][]source}{}{}{}}
\makeatother
\crefname{originalreference}{original reference}{original references}
\Crefname{originalreference}{Original reference}{Original references}
\makeatletter
\newlabel{lemma-locally-acyclic-quasi-finite-base-change}{{0GJS}{source}{}{sourcepointer-2}{}}
\newlabel{lemma-locally-acyclic-quasi-finite-base-change@cref}{{[originalreference][0][]0GJS}{[1][0][]source}{}{}{}}
\makeatother
\crefname{originalreference}{original reference}{original references}
\Crefname{originalreference}{Original reference}{Original references}
\makeatletter
\newlabel{lemma-stalk-pullback}{{008H}{source}{}{sourcepointer-3}{}}
\newlabel{lemma-stalk-pullback@cref}{{[originalreference][0][]008H}{[1][0][]source}{}{}{}}
\makeatother
\crefname{originalreference}{original reference}{original references}
\Crefname{originalreference}{Original reference}{Original references}
\makeatletter
\newlabel{theorem-exactness-stalks}{{03PU}{source}{}{sourcepointer-4}{}}
\newlabel{theorem-exactness-stalks@cref}{{[originalreference][0][]03PU}{[1][0][]source}{}{}{}}
\makeatother
\crefname{originalreference}{original reference}{original references}
\Crefname{originalreference}{Original reference}{Original references}
\makeatletter
\newlabel{theorem-higher-direct-images}{{03Q9}{source}{}{sourcepointer-5}{}}
\newlabel{theorem-higher-direct-images@cref}{{[originalreference][0][]03Q9}{[1][0][]source}{}{}{}}
\makeatother
\title{The comparison map beta between the stalk of K and the cohomology of the Milnor fibre is an isomorphism}
\author{The Stacks Project Authors}
\date{stacks-project:etale-cohomology.tex:local-snapshot-20261011; GNU FDL 1.2 or later}
\begin{document}
\maketitle
\noindent\textit{Editorial note.} All delivered mathematical text is the author's own, copied from the pinned Stacks Project chapter file etale-cohomology.tex; no statement or proof step was rewritten, reordered or removed.\par\medskip
\noindent\textit{Editorial note.} The counted proof is the author's complete argument: the proof is written out in full: the reduction to the strictly henselian case with the observation that local acyclicity is preserved, the description of the relevant fibre product, the identification of the stalk map with alpha, and the concluding stalk criterion.\par\medskip
\noindent\textit{Editorial note.} The label lemma-beta-is-isomorphism names the counted statement, and the proof environment immediately following it in the source is delivered verbatim as its proof.\par\medskip
\noindent\textit{Editorial note.} The counted statement is carried by the author's own lemma environment in the chapter source.\par\medskip
\noindent\textit{Editorial note.} The statement refers to maps constructed earlier in the chapter and is delivered verbatim; the five cited inputs are kept as external pointers with their canonical Stacks tags.\par\medskip
\noindent\textit{Editorial note.} The preamble is the chapter's own native preamble excerpt already used for this source: the Xy-pic setup (native preamble lines 14-18) and the native theorem declarations with the mathematical macros (native preamble lines 154-217).\par\medskip
\noindent\textit{Allowed standard prerequisites (editorial).}\begin{itemize}
\item Local acyclicity of f relative to K via the maps alpha and beta. Same chapter, the definition of local acyclicity using the maps alpha\_\{K,x,t\} and beta\_\{K,s,t\} and their relation. Supplies the map beta whose isomorphism is asserted and the description of alpha that it is compared with.
\item Stalks and pullback: i\textasciicircum{}\{-1\} preserves stalks. Same chapter, the lemma identifying the stalk of a pullback with the stalk of the original sheaf. Supplies the first equality in the computation of the stalk of i\textasciicircum{}\{-1\}beta.
\item Cohomology of the Milnor fibre via higher direct images. Same chapter, the theorem describing the stalk of Rh\_* as the cohomology of the fibre F\_\{x,t\}. Supplies the identification of Rh\_*(K|\_\{X\_t\})\_\{x\} with R Gamma(F\_\{x,t\}, K).
\item A map of sheaves is an isomorphism iff it is on stalks. Same chapter, the exactness-on-stalks theorem for the etale topology. Supplies the final reduction from sheaves to stalks.
\end{itemize}\medskip
\ifdefined\hypertarget\hypertarget{sourcepointer-1}{}\fi
\noindent\textit{Source reference.} Original reference 0GJP: sources/stacks-etale-cohomology/pinned-original-source.tex, source lines 19110--19110. Prerequisite pointer invoked at proof line(s) 19459; it names a standard fact of the same chapter that the counted argument uses, and every step of the counted proof built on it is retained verbatim in the delivered body.\par\medskip
\ifdefined\hypertarget\hypertarget{sourcepointer-2}{}\fi
\noindent\textit{Source reference.} Original reference 0GJS: sources/stacks-etale-cohomology/pinned-original-source.tex, source lines 19289--19289. Prerequisite pointer invoked at proof line(s) 19431; it names a standard fact of the same chapter that the counted argument uses, and every step of the counted proof built on it is retained verbatim in the delivered body.\par\medskip
\ifdefined\hypertarget\hypertarget{sourcepointer-3}{}\fi
\noindent\textit{Source reference.} Original reference 008H: sources/stacks-etale-cohomology/pinned-original-source.tex, source lines 4649--4649. Prerequisite pointer invoked at proof line(s) 19444; it names a standard fact of the same chapter that the counted argument uses, and every step of the counted proof built on it is retained verbatim in the delivered body.\par\medskip
\ifdefined\hypertarget\hypertarget{sourcepointer-4}{}\fi
\noindent\textit{Source reference.} Original reference 03PU: sources/stacks-etale-cohomology/pinned-original-source.tex, source lines 3546--3546. Prerequisite pointer invoked at proof line(s) 19461; it names a standard fact of the same chapter that the counted argument uses, and every step of the counted proof built on it is retained verbatim in the delivered body.\par\medskip
\ifdefined\hypertarget\hypertarget{sourcepointer-5}{}\fi
\noindent\textit{Source reference.} Original reference 03Q9: sources/stacks-etale-cohomology/pinned-original-source.tex, source lines 8015--8015. Prerequisite pointer invoked at proof line(s) 19456; it names a standard fact of the same chapter that the counted argument uses, and every step of the counted proof built on it is retained verbatim in the delivered body.\par\medskip
\begin{lemma}
\label{lemma-beta-is-isomorphism}
The map $i^{-1}\beta_{K, \overline{s}, \overline{t}}$ is an isomorphism.
\end{lemma}

\begin{proof}
The construction of the maps $h$, $i$, $\beta_{K, \overline{s}, \overline{t}}$
only depends on the
base change of $X$ and $K$ to $\Spec(\mathcal{O}^{sh}_{S, \overline{s}})$.
Thus we may and do assume that $S$ is a strictly henselian scheme
with closed point $\overline{s}$. Observe that the local acyclicity
of $f$ relative to $K$ is preserved by this base change (for example
by Lemma \ref{lemma-locally-acyclic-quasi-finite-base-change} or just
directly by comparing strictly henselian rings in this very special case).

\medskip\noindent
Let $\overline{x}$ be a geometric point of $X_{\overline{s}}$.
Or equivalently, let $\overline{x}$ be a geometric point whose image
by $f$ is $\overline{s}$. Let us compute the stalk of
$i^{-1}\beta_{K, \overline{s}, \overline{t}}$
at $\overline{x}$. First, we have
$$
(i^{-1}\beta_{K, \overline{s}, \overline{t}})_{\overline{x}} =
(\beta_{K, \overline{s}, \overline{t}})_{\overline{x}}
$$
since pullback preserves stalks, see Lemma \ref{lemma-stalk-pullback}.
Since we are in the situation $S = \Spec(\mathcal{O}^{sh}_{S, \overline{s}})$
we see that $h : X_{\overline{t}} \to X$ has the property that
$X_{\overline{t}} \times_X \Spec(\mathcal{O}^{sh}_{X, \overline{x}})
= F_{\overline{x}, \overline{t}}$. Thus we see that
$$
(\beta_{K, \overline{s}, \overline{t}})_{\overline{x}} :
K_{\overline{x}}
\longrightarrow
Rh_*(K|_{X_{\overline{t}}})_{\overline{x}} =
R\Gamma(F_{\overline{x}, \overline{t}}, K)
$$
where the equal sign is Theorem \ref{theorem-higher-direct-images}.
It follows that the map $(\beta_{K, \overline{s}, \overline{t}})_{\overline{x}}$
is none other than the map $\alpha_{K, \overline{x}, \overline{t}}$ used
in Definition \ref{definition-locally-acyclic}. The result follows as we may
check whether a map is an isomorphism in stalks by
Theorem \ref{theorem-exactness-stalks}.
\end{proof}

\end{document}
